\documentclass[11pt]{article}
\usepackage[a4paper,top=0.7in,bottom=0.7in,left=0.9in,right=0.9in]{geometry}
\usepackage{parskip}
\usepackage{times}
\usepackage[hidelinks]{hyperref}
\usepackage{microtype}
\usepackage{enumitem}
\pagestyle{empty}
\setlength{\parskip}{3pt plus 1pt}
\setlength{\emergencystretch}{2em}

\begin{document}

{\raggedleft\small
Md Habibur Rahman (corresponding author) \\
Department of AI Convergence Engineering, Gyeongsang National University \\
Jinju, South Korea \quad\textbullet\quad \href{mailto:habib@gnu.ac.kr}{habib@gnu.ac.kr} \\[0.4em]
\par}

\noindent The Editorial Office, \textit{Quantum Reports} (MDPI)

\vspace{0.5em}

\noindent\textbf{Re: Submission of ``Predicting the Trainability of Variational
Quantum Circuits: A Data-Driven Model for Barren Plateaus''}

\vspace{0.3em}

\noindent Dear Editors,

I am pleased to submit my manuscript, \textit{``Predicting the Trainability of
Variational Quantum Circuits: A Data-Driven Model for Barren Plateaus,''} for
consideration as an original research article in \textit{Quantum Reports}.

\textbf{Significance.} Barren plateaus---the exponential vanishing of
cost-function gradients as a variational quantum circuit is scaled---are widely
regarded as the central obstacle to training quantum machine-learning models.
Diagnosing them today requires expensive gradient sampling on each candidate
circuit, infeasible precisely in the large-qubit regime where the problem matters
most. This manuscript reframes that diagnosis as a \emph{supervised prediction}
problem: can a purely classical model predict a circuit's trainability from its
architecture alone---qubit count, depth, ansatz type, entanglement pattern,
entangler gate, and cost-observable locality---without evaluating a single
gradient? Using an exact statevector simulator I generate a labeled dataset of
20{,}000 random circuits and show that a gradient-boosted model predicts the
gradient-variance trainability with high accuracy (AUC 0.99 for the
barren/trainable decision) and, most notably, \emph{extrapolates} from small
circuits to the largest sizes that can still be simulated---yielding a cheap
ansatz-screening heuristic and an interpretable account of which design choices
drive untrainability.

\textbf{Context and relationship to existing work.} The barren-plateau literature
is overwhelmingly analytical, deriving variance scaling laws for particular ansatz
families (e.g., McClean et al.\ 2018; Cerezo et al.\ 2021; Larocca et al.\ 2025).
This contribution is complementary and empirical: I learn a predictive map across a
heterogeneous design space and show that it recovers the established theory from
data. Consistent with recent skeptical benchmarking (Bowles et al.\ 2024;
Schuld \& Killoran 2022), I make \textbf{no} claim of quantum advantage---the
object predicted is a property \emph{of} quantum circuits, while the predictor is
a classical diagnostic tool.

\textbf{Fit to the scope of Quantum Reports.} The manuscript sits squarely within
the journal's remit of quantum computing, quantum information, and quantum machine
learning, and is a natural fit for the Special Issue \textit{``The Integration of
Quantum Computing with Artificial Intelligence.''}

\textbf{Reproducibility.} In keeping with the journal's open-science standards, all
code, fixed random seeds, and the full dataset are openly available at
\url{https://github.com/bithabib/quantum_computer_simulator}, and the underlying
simulator runs interactively in the browser at
\url{https://quantum.qbithabib.com}; every reported number regenerates from a
single command.

\textbf{Required statements.}
\begin{itemize}[nosep,leftmargin=1.4em,itemsep=1.5pt]
\item I confirm that neither the manuscript nor any parts of its content are
currently under consideration for publication with or published in another journal.
\item I am the sole author of this manuscript and approve its submission to
\textit{Quantum Reports}.
\item This manuscript has not been previously submitted to any MDPI journal.
\item Generative AI (Claude, Anthropic) was used for language editing and grammar
correction only; this is disclosed in the Acknowledgments of the manuscript.
\item The author declares no conflicts of interest.
\end{itemize}

I believe this work will interest the readership of \textit{Quantum Reports} and
look forward to the reviewers' feedback.

\vspace{0.3em}
\noindent Sincerely, \\[0.6em]
Md Habibur Rahman \\
Corresponding author --- \href{mailto:habib@gnu.ac.kr}{habib@gnu.ac.kr} \\
Department of AI Convergence Engineering, Gyeongsang National University, Jinju,
South Korea

\end{document}
